\chapter{Iron Deficiency}
\index{Iron Deficiency}

\section*{Definition and Overview}
Iron deficiency is a lack of iron in the body, occurring when intake is inadequate, absorption is impaired, or losses are increased. Because iron is needed to make haemoglobin, the oxygen-carrying protein in red blood cells, deficiency causes anaemia, with fatigue, breathlessness, and paleness. Iron deficiency is the most common nutritional deficiency in the world and is particularly common in women of reproductive age, who lose iron through menstruation, and in pregnancy, when demands rise sharply. It also affects babies and toddlers, teenagers, vegetarians and vegans, and people with gut conditions or chronic bleeding. Diagnosis is simple with blood tests, and treatment with dietary change and iron supplements is highly effective.

\section*{Epidemiology}
Iron deficiency affects around one in eight women of reproductive age and a similar proportion of teenage girls, and it is the most common cause of anaemia worldwide, affecting around a quarter of the global population. Men and post-menopausal women are less often affected, and deficiency in them usually indicates blood loss from the gut, requiring investigation. Pregnancy increases iron needs markedly, and iron deficiency affects up to half of pregnancies. Infants and toddlers are at risk when iron stores and breast milk alone do not meet needs. People with coeliac disease, inflammatory bowel disease, and after gastric surgery have impaired absorption.

\section*{Aetiology and Pathophysiology}
Iron is absorbed mainly in the upper small intestine and is used to build haemoglobin in red blood cells, myoglobin in muscle, and enzymes throughout the body. The body has no effective way to excrete iron, so balance is maintained by regulating absorption. Deficiency develops when intake is low, as in vegetarian and vegan diets that lack haem iron; when absorption is impaired by coeliac disease, gastritis, or surgery; or when losses exceed intake, through menstruation, pregnancy, repeated blood donation, or gut bleeding. As stores, measured by ferritin, are depleted, haemoglobin production falls and anaemia develops. Iron deficiency without anaemia can still cause symptoms, including fatigue and impaired concentration.

\section*{Clinical Features}
\begin{itemize}
    \item Fatigue, weakness, and reduced exercise tolerance
    \item Breathlessness on exertion and palpitations
    \item Paleness of the skin, nail beds, and inner eyelids
    \item Headaches, dizziness, and difficulty concentrating
    \item Cold hands and feet
    \item Pica: cravings for ice, dirt, or other non-food items
    \item Brittle or spoon-shaped nails and hair thinning
    \item Sore tongue and cracks at the corners of the mouth
    \item Restless legs, which is linked to iron deficiency
    \item In children: poor growth, irritability, and impaired learning
\end{itemize}

\section*{Differential Diagnosis}
Anaemia and fatigue have other causes that must be distinguished.
\begin{itemize}
    \item \textbf{Other nutritional deficiencies:} vitamin B12 and folate deficiency
    \item \textbf{Anaemia of chronic disease:} from inflammation, infection, or cancer
    \item \textbf{Thalassaemia and haemoglobinopathies:} inherited blood conditions
    \item \textbf{Chronic kidney disease:} reduced erythropoietin production
    \item \textbf{Bleeding:} gut, menstrual, or surgical blood loss
    \item \textbf{Hypothyroidism, depression, and chronic fatigue:} other causes of tiredness
    \item \textbf{Bone marrow disorders:} rarer causes of anaemia
\end{itemize}

\section*{When to Seek Medical Care}
Anyone with persistent fatigue, breathlessness, or paleness should have a blood test. Women with heavy periods, people with gut symptoms, vegetarians and vegans, and anyone with unexplained blood loss should be assessed. Men and post-menopausal women with iron deficiency need investigation of the gut to exclude bleeding. Pregnant women are screened routinely and should take iron as advised.

\textbf{Contact your local emergency services for:}
\begin{itemize}
    \item Severe breathlessness, chest pain, or collapse
    \item Black, tarry, or bloody stools
    \item Vomiting blood
    \item Heavy bleeding with faintness or dizziness
    \item Sudden severe weakness or confusion
\end{itemize}

\section*{Diagnosis and Investigations}
\begin{itemize}
    \item \textbf{Full blood count:} haemoglobin and red cell indices
    \item \textbf{Ferritin:} the best test of iron stores, with low levels indicating depletion
    \item \textbf{Iron studies:} serum iron, transferrin, and transferrin saturation
    \item \textbf{B12 and folate:} to exclude other deficiencies
    \item \textbf{Coeliac screening:} in people with gut symptoms or unexplained deficiency
    \item \textbf{Investigation of bleeding:} endoscopy or colonoscopy in men and post-menopausal women
    \item \textbf{Gynaecological review:} for heavy menstrual bleeding in women
    \item \textbf{Pregnancy screening:} routine iron and haemoglobin checks
\end{itemize}

\section*{Management}
\begin{itemize}
    \item \textbf{Oral iron supplements:} iron tablets, taken with vitamin C to aid absorption and away from tea and coffee
    \item \textbf{Dietary change:} increasing iron-rich foods, including red meat, and enhancing absorption with vitamin C
    \item \textbf{Treatment of the cause:} treating coeliac disease, gut bleeding, or heavy periods
    \item \textbf{Intravenous iron:} for people who cannot tolerate or absorb oral iron, and in late pregnancy
    \item \textbf{Blood transfusion:} for severe or symptomatic anaemia
    \item \textbf{Monitoring:} repeat blood tests to confirm response
    \item \textbf{Continuing supplementation:} to replenish stores after the anaemia resolves
    \item \textbf{Baby and child care:} iron-fortified foods and screening in at-risk groups
\end{itemize}

\section*{Complications}
\begin{itemize}
    \item Chronic fatigue and reduced quality of life
    \item Impaired exercise performance and work productivity
    \item Complications in pregnancy: preterm birth, low birth weight, and maternal fatigue
    \item Impaired cognitive development in infants and children
    \item Heart strain from long-standing anaemia
    \item Restless legs and its effect on sleep
    \item Missed underlying disease, such as bowel cancer, when bleeding is not investigated
\end{itemize}

\section*{Prognosis and Outcomes}
Iron deficiency responds well to treatment. Symptoms usually improve within weeks of starting supplements, and haemoglobin returns to normal within one to two months, with iron stores replenished over further months. When the underlying cause is treated, recurrence is uncommon. In pregnancy, early treatment protects both mother and baby. The most important step is identifying the cause, because iron deficiency in men and older women may be the first sign of serious gut disease.

\section*{Prevention and Self-Care}
\begin{itemize}
    \item Eat iron-rich foods: red meat, poultry, fish, legumes, eggs, and fortified cereals
    \item Combine iron-rich plant foods with vitamin C, such as fruit or vegetables
    \item Avoid tea, coffee, and dairy at the same time as iron-rich meals
    \item Take iron supplements as prescribed, with water and vitamin C
    \item Take iron during pregnancy as advised
    \item Treat heavy periods and seek gynaecological review
    \item Have blood tests for persistent fatigue
    \item Feed babies iron-fortified foods from six months
    \item Consider a dietitian review for vegetarian and vegan diets
\end{itemize}

\section*{Sources and Further Reading}
The following reputable sources provide further information for healthcare students and patients seeking reliable, current advice about iron deficiency.
\begin{itemize}
    \item Healthdirect Australia. \textit{Iron deficiency.} https://www.healthdirect.gov.au/
    \item The Royal Australian College of General Practitioners. \textit{Iron deficiency guideline.} https://www.racgp.org.au/
    \item Haematology Society of Australia and New Zealand. \textit{Iron deficiency management.} https://www.hsanz.org.au/
    \item Dietitians Australia. \textit{Iron in the diet.} https://dietitiansaustralia.org.au/
    \item Pregnancy, Birth and Baby. \textit{Iron in pregnancy.} https://www.pregnancybirthbaby.org.au/
    \item NHS. \textit{Iron deficiency anaemia.} https://www.nhs.uk/conditions/iron-deficiency-anaemia/
\end{itemize}
